\documentclass[12pt,a4paper]{article}
\usepackage[utf8]{inputenc}
\usepackage[italian]{babel}
\usepackage{amsmath, amssymb, amsfonts, bm}
\usepackage{graphicx}
\usepackage{pgfplots}
\pgfplotsset{compat=1.18}
\usepackage{geometry}
\geometry{margin=2.5cm}
\usepackage{caption}
\usepackage{hyperref}

\title{\centering Studio matematico sulla coscienza umana \\ Legge 3: Legge della Volontà}
\author{Autore: Giuseppe Carlino}
\date{\today}

\begin{document}

\maketitle

\section{Premessa}

La \textbf{Legge della Volontà} descrive come la coscienza umana trasformi la conoscenza accumulata e le scelte effettuate in azioni volontarie.  
In questo modello, la volontà $\bm{w}$ non è un processo indipendente: è il risultato della combinazione integrata dei tre stati precedenti:

\[
\bm{w} = \bm{w}(t_c, \bm{k}, s, F),
\]

dove:

\begin{itemize}
    \item $t_c$ è il tempo percepito (Legge 1),
    \item $\bm{k}$ è la conoscenza accumulata (Legge 0),
    \item $s$ è la scelta effettuata (Legge 2),
    \item $F = \{f_1, f_2, \dots, f_n\}$ rappresenta fattori esterni o interni che influenzano la complessità del comportamento volontario.
\end{itemize}

Il principio fondamentale è che la volontà umana emerge dall’interazione dei processi cognitivi, decisionali e temporali.

\section{Formalizzazione concettuale}

La volontà $\bm{w}$ è una variabile continua:

\[
\bm{w} \in [0, \bm{w}_{\max}],
\]

con $\bm{w}_{\max}$ massimo livello di intensità volontaria.

Una funzione che integra i contributi dei tre stati precedenti:

\begin{equation}
\bm{w}(t) = \phi(t_c, \bm{k}(t), s(t), F, \bm{\xi}_w(t)),
\end{equation}

dove $\bm{\xi}_w(t)$ rappresenta fluttuazioni interne o perturbazioni.

Una forma esplicita:

\begin{equation}
\bm{w}(t) = \gamma \, s(t) \, \frac{\|\bm{k}(t)\|}{\|\bm{k}_{\max}\|} \left(1 + \sum_{i=1}^{n} \alpha_i f_i \right) + \bm{\xi}_w(t),
\end{equation}

con:

\begin{itemize}
    \item $\gamma>0$ coefficiente di trasformazione della scelta in volontà,
    \item $\alpha_i$ pesi relativi ai fattori esterni $f_i$,
    \item $\bm{\xi}_w(t)$ perturbazioni stocastiche.
\end{itemize}

\subsection{Forma dinamica}

L’evoluzione della volontà può essere descritta da:

\begin{equation}
\frac{d\bm{w}}{dt} = -\nu (\bm{w} - \bm{w}_e(\bm{k},s,F)) + \bm{\eta}_w(t),
\end{equation}

dove:

\begin{itemize}
    \item $\bm{w}_e(\bm{k},s,F) = \gamma s \frac{\|\bm{k}\|}{\|\bm{k}_{\max}\|} (1 + \sum \alpha_i f_i)$ equilibrio volontà-istante,
    \item $\nu>0$ coefficiente di richiamo verso $\bm{w}_e$,
    \item $\bm{\eta}_w(t)$ perturbazioni interne o esterne.
\end{itemize}

---

\section{Diagramma concettuale 3D}

\begin{center}
\begin{tikzpicture}
\begin{axis}[
    view={60}{30},
    width=14cm, height=10cm,
    axis lines=center,
    grid=major,
    xlabel={$s$},
    ylabel={$\|\bm{k}\|$},
    zlabel={$\bm{w}$},
    xmin=0, xmax=10,
    ymin=0, ymax=5,
    zmin=0, zmax=5,
    colormap/viridis
]

% Traiettoria concettuale della volontà
\addplot3[smooth, thick, purple] coordinates {
(0,0,0.2) (1,0.5,0.5) (2,1,1.0) (3,1.5,1.5) (4,2,2.0)
(5,2.5,2.5) (6,3,3.0) (7,3.5,3.5) (8,4,4.0) (9,4.5,4.2)
};

% Nuvola concettuale di fattori esterni F
\addplot3[only marks, mark=*, mark size=0.5pt, gray, opacity=0.3] coordinates {
(5,2.5,2.5) (5.2,2.6,2.7) (4.8,2.4,2.3) (5.1,2.55,2.6) (4.9,2.45,2.4)
};

\end{axis}
\end{tikzpicture}
\end{center}

\caption*{Lo spazio tridimensionale rappresenta la volontà $\bm{w}$ come funzione di conoscenza $\|\bm{k}\|$, scelta $s$ e fattori esterni $F$. La traiettoria viola mostra l’evoluzione concettuale della volontà verso l’equilibrio, la nuvola grigia rappresenta fluttuazioni esterne.}

---

\section{Interpretazione e implicazioni}

\begin{enumerate}
    \item La volontà integra tempo, conoscenza e scelta.
    \item Le perturbazioni $\bm{\eta}_w$ rappresentano incertezza e variabilità.
    \item Sistemi artificiali possono simulare scelte, ma non possiedono retroazione continua tra conoscenza, scelta e volontà.
    \item La Legge della Volontà prepara la formalizzazione generale della \textbf{Legge Quadristato}.
    \item L’uomo conserva capacità decisionale e volontaria non replicabile da algoritmi predeterminati.
\end{enumerate}

\section{Relazione con le altre leggi}

\begin{itemize}
    \item Legge 0 (Conoscenza) fornisce $\bm{k}(t)$ come base cognitiva.
    \item Legge 1 (Tempo) regola $t_c$, influenzando conoscenza e volontà.
    \item Legge 2 (Scelta) fornisce $s(t)$ come input fondamentale per la volontà.
    \item Legge Generale (Quadristato) unifica $t_c$, $\bm{k}$, $s$ e $\bm{w}$ in un modello coerente della coscienza complessa.
\end{itemize}

\end{document}